% Source: https://artofproblemsolving.com/community/c3h3062857_1000_post_special_math_story_and_some_stuff
% Stable identity: aops-topic-3062857-q01
(a) The pyramid apices are $(\pm2,0,0),(0,\pm2,0),(0,0,\pm2)$. Every edge of the resulting rhombic dodecahedron joins such an apex to a cube corner at distance $\sqrt{1+1+1}=\sqrt3$. Adjacent triangular pyramid faces sharing a cube edge are coplanar and form one rhombic face. The volume is $8+6\left(\dfrac{1}{3}\cdot4\cdot1\right)=\boxed{16}$. (b) Scaling the edge from $\sqrt3$ to $n$ gives volume $16\left(\dfrac{n}{\sqrt3}\right)^3=\dfrac{16n^3}{3\sqrt3}$. The determinant of the three displayed vectors has absolute value $\dfrac{4}{3\sqrt3}$, which is the volume of one parallelepiped. Dividing volumes shows that any such tiling contains $\boxed{4n^3}$ tiles. This is a necessary tile count under the stated tiling assumption, not a proof of tiling existence.
\begin{asy}
usepackage("amsmath");
size(7cm);
pair proj(real x,real y,real z){return (x+.45*y,z+.3*y);}pair[] corners;for(int x=-1;x<=1;x+=2)for(int y=-1;y<=1;y+=2)for(int z=-1;z<=1;z+=2){pair p=proj(x,y,z);draw(p--proj(2*x,0,0),gray);draw(p--proj(0,2*y,0),gray);draw(p--proj(0,0,2*z),gray);}dot(proj(0,0,2));dot(proj(2,0,0));dot(proj(0,2,0));label("$(0,0,2)$",proj(0,0,2),N);label("$(2,0,0)$",proj(2,0,0),E);label("$(0,2,0)$",proj(0,2,0),NE);label("$(1,1,1)$",proj(1,1,1),NE);label("rhombic dodecahedron",(0,-2.5));
\end{asy}
